\documentclass{situationist-preprint}

\title{Presentism and the Truthmaker Objection}
\msid{TSM-03-b}
\msfamily{TSM-03}
\msdiscipline{Time, Space and Modality}
\mstype{Companion manuscript}
\mskeywords{presentism; truthmaking; Lucretianism; ersatz times; cheating; past}

\begin{document}
\maketitle

\begin{abstract}
The truthmaker objection says that presentism cannot account for what makes truths about the past true. Presentists have replied with past-directed properties of the present world, with ersatz times, and with the denial that truth needs grounding. TSM-03-a uses the objection in a weak, supervenience form. We show that the weak form is the right one to press: it is accepted by the presentist's own replies, it does not depend on contested truthmaker maximalism, and every reply that satisfies it does so by building the past into a present fact. The objection therefore does not merely embarrass presentism; it shows that presentism's surrogates are the past under another description.
\end{abstract}

\section{Two forms of the objection}

The strong form of the objection assumes truthmaker maximalism: every truth has a truthmaker, an entity whose existence necessitates it \citep{Armstrong2004}. The presentist need not accept maximalism, and many truthmaker theorists do not; the strong form invites the reply that negative existentials and past truths alike are exceptions. The weak form is Principle~P1 of TSM-03-a: truth supervenes on what exists and on which fundamental properties and relations existing things instantiate \citep{Bigelow1988,Lewis2001}. It requires no truthmaker for any particular truth, only that two situations alike in being be alike in truth.

\begin{proposition}\label{p:accept}
Each of the main presentist replies accepts the weak form.
\end{proposition}

\begin{proof}
Lucretianism posits past-tensed properties of the present world, such as having been such that Caesar crossed the Rubicon, precisely in order that past truths supervene on how the present is \citep{Bigelow1996}. Ersatzist presentism posits abstract, presently existing representations of past times and a present relation to them, for the same purpose \citep{Crisp2007}. Each is offered as a way of meeting the supervenience requirement, not of escaping it.
\end{proof}

Only the third reply, that truths about the past need no grounding at all \citep{Merricks2007}, rejects the weak form, and TSM-03-a answers it separately (Objection~2 there). The remaining question is whether the first two succeed.

\section{The surrogates are the past}

Lemma~3.1 of TSM-03-a argues that they do not, because their properties are defined by what is true at other times. We make the point more general.

\begin{proposition}\label{p:dilemma}
Let $F$ be a present feature offered to ground a truth $p$ about the past, in a case covered by Assumption~2.3 of TSM-03-a, where the intrinsic present and the laws do not entail $p$. Then either $F$ is intrinsic to what is present, and it does not ground $p$, or $F$ is specified as the property of being such that $p$ was the case, and its instantiation is the fact that $p$ was the case.
\end{proposition}

\begin{proof}
If $F$ is intrinsic, then since the intrinsic present and the laws do not entail $p$, there is a situation alike in $F$ in which $p$ is false, and $F$ does not make the difference. If $F$ is not intrinsic, it is individuated by its relation to how things were. A Lucretian property is so individuated by its specification, and the presently existing ersatz time that \emph{was} accurate is so individuated by the past-tensed fact of its accuracy. In either case, that $F$ is instantiated is the fact that $p$ was so.
\end{proof}

Proposition~\ref{p:dilemma} explains the charge that such views \enquote{cheat} \citep{Sider2001}. The cheating is not that the posited properties are unfamiliar. It is that they are introduced as present but have, as their whole content, the past fact. A world with such properties is one in which the past is stored in the present by stipulation, and the storage has no features beyond the fact stored. That is a difference of description, not of ontology: what the presentist calls the world's having been such that the coin landed heads, the eternalist calls the coin's landing heads at an earlier time.

\section{Objection and reply}

\begin{objection}
Presentism can at least explain why the past is fixed and the future open: past-tensed properties are instantiated and future-tensed ones are not.
\end{objection}

\begin{reply}
The asymmetry is better explained without presentism. In a branching structure, relative to any moment there is one past and several possible continuations (TSM-12); the past is fixed because it is a single region, the future open because it is not. TSM-03-a states its thesis in these terms: for every moment, its past exists, and no single future region is asserted to exist.
\end{reply}

\section{Conclusion}

The truthmaker objection should be pressed in its weak form, which the presentist's own replies accept. The replies that satisfy it do so by positing present features whose content is the past fact; they relocate the past rather than replace it. The best version of presentism thus concedes the substance of TSM-03: the past is real, and is elsewhere in time.

\begin{thebibliography}{9}
\bibitem[Armstrong(2004)]{Armstrong2004} Armstrong, D.~M. (2004). \emph{Truth and Truthmakers}. Cambridge University Press.
\bibitem[Bigelow(1988)]{Bigelow1988} Bigelow, J. (1988). \emph{The Reality of Numbers: A Physicalist's Philosophy of Mathematics}. Clarendon Press.
\bibitem[Bigelow(1996)]{Bigelow1996} Bigelow, J. (1996). Presentism and properties. \emph{Philosophical Perspectives}, 10, 35--52.
\bibitem[Crisp(2007)]{Crisp2007} Crisp, T.~M. (2007). Presentism and the grounding objection. \emph{No\^us}, 41(1), 90--109.
\bibitem[Lewis(2001)]{Lewis2001} Lewis, D. (2001). Truthmaking and difference-making. \emph{No\^us}, 35(4), 602--615.
\bibitem[Merricks(2007)]{Merricks2007} Merricks, T. (2007). \emph{Truth and Ontology}. Oxford University Press.
\bibitem[Sider(2001)]{Sider2001} Sider, T. (2001). \emph{Four-Dimensionalism: An Ontology of Persistence and Time}. Oxford University Press.
\end{thebibliography}

\end{document}
